\SourcePage{327}
\SGAsection{}{原书术语索引}
\noindent 下列定位保留原书的讲次与条目编号；词头按原书次序排列。
\begin{longtable}{@{}p{0.75\textwidth}p{0.19\textwidth}@{}}
可达（函子；Accessible） & I 9.2\\
可达（对象；Accessible） & I 9.3\\
子拓扑斯的闭包（Adhérence） & IV 9.4.8\\
阿廷（集合；Artinien） & I $0$, I 11.6\\
阿廷（对象；Artinien） & I 8.12.6\\[3pt]
双覆盖（态射、族；Bicouvrant） & II 5.2\\
笛卡儿、余笛卡儿（Cartésien, cocartésien） & I 10.1\\
$\mathcal U$-范畴 & I 1.1\\
滤过范畴（Catégorie filtrante） & I 2.7\\
（拓扑斯的）点范畴 & IV 4.6.3\\
伪滤过范畴（Catégorie pseudo-filtrante） & I 2.7\\
共尾（函子、子范畴；Cofinal） & I 8.1.1\\
余生成（子范畴；Cogénératrice） & I 7.1\\
比较引理（Comparaison） & III 4.1\\
常值、本质常值（Ind-对象；Constant, essentiellement constant） & I 8.4\\
覆盖（筛；Couvrant） & II 1.1.1\\
覆盖（态射、族；Couvrant） & II 5.2\\
覆盖族（Couvrante） & II 1.2\\
筛（Crible） & I 4.1\\[3pt]
支配（拓扑斯态射；Dominant） & IV 8.8\\
泛有效满态射族（Épimorphique effective universelle） & II 2.5\\
满态射族（Épimorphique） & I 10.3\\
泛严格满态射（筛、族；Épimorphique strict universel） & II 2.5\\
泛满态射族（Épimorphique universelle） & I 10.3\\
满态射、严格满态射（Épimorphisme, épimorphisme strict） & I 10.3\\
\end{longtable}
\SourcePage{328}
\begin{longtable}{@{}p{0.75\textwidth}p{0.19\textwidth}@{}}
有效满态射（Épimorphisme effectif） & I 10.3\\
本质（态射、点；Essentiel） & IV 7.6\\
展拓扑斯（Étendue） & IV 9.8.2\\
子拓扑斯的外部（Extérieur） & IV 9.4.8\\[3pt]
（取值于范畴 $D$ 的）层 & II 6.1\\
（集合）层 & II 2.1\\
态射层 & IV 10.2\\
粘合层 & IV 9.6\\
基数滤过（Filtration cardinale） & I 9.12\\
余连续函子 & III 2.1\\
保守函子 & I 6.1\\
连续函子 & III 1.1\\
粘合函子 & IV 9.3.5\\
左正合（右正合、正合）函子 & I 2.4\\
纤维函子 & IV 4.6.3\\
忠实函子 & I 6.1\\
零延拓函子 & IV 11.3.1\\
限制标量函子 & IV 11.1.4\\
截面函子 & IV 4.3.6\\
纤维函子 & IV 6.0\\
子拓扑斯的边界 & IV 9.4.8\\[3pt]
生成、余生成（族） & I 7.1\\
生成、余生成（子范畴） & I 7.1, I 7.9\\
（景的）生成族 & II 3.0.1\\
（景的）拓扑生成族 & II 3.0.1\\
（点的）一般化（Générisation） & IV 4.2.2\\
（拓扑空间的）大拓扑斯（Gros topos） & IV 4.10, IV 2.5\\[3pt]
模的正像 & IV 13.2.1\\
拓扑斯态射的像 & IV 9.1.7\\
（子拓扑斯的）逆像 & IV 9.1.6\\
（诱导拓扑斯的）逆像 & IV 5.10\\
模的逆像 & IV 13.2.2\\
不可达（基数） & I $0$\\
不可达（基数） & I 11.5\\
包含态射（Inclusion） & IV 5.2\\
Ind-伴随 & I 8.11.1\\
Ind-对象 & I 8\\
\end{longtable}
\SourcePage{329}
\begin{longtable}{@{}p{0.75\textwidth}p{0.19\textwidth}@{}}
严格 Ind-对象 & I 8.12.1\\
Ind-可表（函子） & I 8.2\\
诱导（拓扑斯；Induit） & IV 5.2\\
子拓扑斯的内部 & IV 9.4.8\\[3pt]
归纳极限 & I 2.3.1\\
泛归纳极限 & I 2.5\\
逆极限 & I 2.1\\
$\mathcal U$-逆极限 & I 2.2.1\\
有限逆极限与归纳极限 & I 2.3.1\\
局部（性质关系；Locale） & IV 8.2\\
局部化函子 & IV 5.4\\
局部化态射 & IV 5.2\\
单态射、单态射（Monomorphisme, monomorphisme） & I 10.4\\
拓扑斯态射 & IV 3.1\\
景态射 & IV 4.9\\
环化拓扑斯态射 & IV 13.1\\
局部环化拓扑斯态射 & IV 13.9\\[3pt]
始对象 & II 4.5\\
（拓扑斯的）开对象 & IV 8.2\\[3pt]
小（集合、群、环、范畴等） & I 1.0\\
$\mathcal U$-小（集合） & I 1.0\\
嵌入 & IV 9.1.1\\
闭嵌入 & IV 9.3.5\\
开嵌入 & IV 9.2.1\\
几何点 & IV 13.9\\
（拓扑斯的）点 & IV 6.0\\
由对象表示的预层 & I 1.3.3\\
可表预层 & I 1.4.1\\
（集合）预层 & I 1.2\\
预拓扑 & II 1.3\\
Pro-伴随 & I 8.11.5\\
Pro-对象 & I 8, I 8.10\\
Pro-可表（函子） & I 8.12\\
Pro-可表（函子） & I 8.10.10\\
投影算子（Projecteur） & I 10.6\\
空延拓 & III 5.3\\
可作基变换（箭头、态射；Quarrable） & I 10.7\\
\end{longtable}
\SourcePage{330}
\begin{longtable}{@{}p{0.75\textwidth}p{0.19\textwidth}@{}}
商、严格有效商、泛商 & I 10.8\\[3pt]
加细 & II 1.1\\
等价关系 & I 10.9\\
局部等价关系 & IV 9.8.3\\
有效等价关系、泛有效等价关系 & I 10.10\\
Weil 限制 & IV 5.2\\
（层的）限制 & IV 5.4, III 5.3\\[3pt]
分离（预层） & II 2.1\\
景 & II 1.1.5\\
$\mathcal U$-景 & II 3.0.2\\
局部赋环景 & IV 13.9\\
清醒（拓扑空间；Sobre） & IV 4.2.1\\
不交和 & II 4.5\\
子对象、严格子对象 & I 10.11\\
子拓扑斯 & IV 9.1.1\\
互补子拓扑斯 & IV 9.1.13\\
有补子拓扑斯 & IV 9.1.13\\
闭子拓扑斯 & IV 9.3.5\\
局部闭子拓扑斯 & IV 9.4.9\\
开子拓扑斯 & IV 9.2.3\\
（点的）特殊化 & IV 4.2.2\\
对下降稳定 & IV 8.2\\
严格 Ind-可表 & I 8.12.1\\
（群、群截面的）支集 & IV 9.3.5\\
$r$-结构 & IV 9.8.3\\
支集、余支集 & IV 8.5\\[3pt]
拓扑 & II 1.1\\
$\mathcal U$-拓扑 & II 3.0.2\\
典范拓扑 & II 2.5\\
离散拓扑 & II 1.1.4\\
粗拓扑与混沌拓扑 & II 1.1.4\\
诱导拓扑 & III 3.1\\
拓扑斯 & IV 1.1\\
$\mathcal U$-拓扑斯 & IV 1.1\\
环化拓扑斯 & IV 11.1.1\\
分类拓扑斯 & IV 2.3, 2.4, 2.5\\
等价的拓扑斯 & IV 3.4\\
\end{longtable}
\SourcePage{331}
\begin{longtable}{@{}p{0.75\textwidth}p{0.19\textwidth}@{}}
有限拓扑斯 & IV 9.1.12\\
始拓扑斯、终拓扑斯 & IV 2.2\\
刚性拓扑斯 & IV 9.8.2 b）\\[3pt]
宇宙 & I $0$\\
宇宙 & I 11.1\\
宇宙公理 & I 11.4\\[3pt]
邻域 & IV 6.8\\
\end{longtable}

\SourcePage{333}
\SGAsection{}{原书符号索引}
\begin{longtable}{@{}p{0.75\textwidth}p{0.19\textwidth}@{}}
$\underline a:C^\wedge\to C^{\sim}$ & II 3.4\\
$\underline aF$ & II 3.5\\
$(B,A,f)$（$f:A\to B$ 为函子） & IV 9.5.1\\
$B_G$ & IV 2.3, 4.5\\
$C^{\sim}$ & II 2.1, 6.3.3\\
$C,C',D,\ldots$，范畴 & I 1.2\\
$C^\wedge,C^\wedge_{\mathcal U}$ & II 6.3.3；IV 2.6, 4.6\\
$C^{\sim}_A,C^{\sim}_{\mathrm{ab}}$ & II 6.3.3\\
$C/F$ & I 3.4.0\\
$\operatorname{Cov}(X)$ & II 1.3\\
$C^{\sim}_{\mathcal U}$ & II 3.6\\
${}_AE,E_A$（$E$ 为拓扑斯） & IV 11.1.2\\
$E_{\mathrm{ab}}$ & IV 11.1.5\\
$E/X$ & IV 5.1\\
$\mathscr{F}\!ib(E)$ & IV 6.2\\
$\operatorname{Filt}(E)$ & I 9.12\\
$\Gamma(E,X),\Gamma(X)$ & IV 4.3.6\\
$\operatorname{Hom}_A(M,N),\operatorname{Hom}(M,N)$ & IV 12.1\\
$\operatorname{Hom}_E(X,Y),\operatorname{Hom}(X,Y)$ & IV 10.2\\
$\mathscr{H}\!om\mathrm{topan}(E,E'),\mathscr{H}\!om\mathrm{toplocan}(E,E')$ & IV 13.9 c）\\
$\mathscr{H}\!om\mathrm{top}(E,E')$ & IV 3.2, 3.3\\
$h(X),h_{\mathcal U}(X)$ & I 1.3\\
$i:C^{\sim}\hookrightarrow C^\wedge$ & \\
$\operatorname{Ind}(C),\operatorname{Ind}_{\mathcal U}(C),\operatorname{Ind}_{\mathcal V}(C,\mathcal U)$ & I 8.2.4.4\\
$\xrightarrow{\sim}$，同构 & \\
$j_{X!}^{\mathrm{ab}},j_{X!}^{\mathrm{ens}}$ & IV 11.3.2\\
\end{longtable}
\SourcePage{334}
\begin{longtable}{@{}p{0.75\textwidth}p{0.19\textwidth}@{}}
$J(X)$ & II 1.1\\
$j_X:E/X\to E$ & IV 5.2\\
$j_{X!},j_X^*,j_{X*}$ & III 11.2.2；IV 5.2.2, 11.2.2, 11.3.1\\
$\mathrm{(Karfiness)}$ & IV 9.1.12\\
$L$（函子） & II 3.1\\
$\operatorname{Lib}$（Women’s Liberation，妇女解放），另见 & I 5.8.2, II 6.5\\
$\displaystyle\lim_I$ & I 2.1\\
$\displaystyle\lim_I$（归纳极限） & \\
$\operatorname{loc}(f)$ & IV 5.5.2\\
$\operatorname{loc}(R),\operatorname{glob}(r)$ & IV 9.8.3 b）\\
$M\otimes_A N$ & IV 12.7\\
$\hookrightarrow$，单态射或函子 & \\
$\operatorname{Morsite}(C',C)$ & IV 4.9.3\\
$\operatorname{Ouv}(X)$ & IV 2.1\\
$\varnothing_{\mathrm{top}}$ & IV 4.4\\
$\Phi_v,\psi_v,\Xi_v,\Lambda_v$ & IV 10.5\\
$\mathscr{P}\!oints(E)$ & IV 4.6.3, 6.1\\
$\mathscr{P}\!t(E)$ & IV 6.1\\
$\mathscr{P}\!t\mathrm{gom}(E)$ & IV 13.9 f）\\
$\operatorname{Pro}(C),\operatorname{Pro}_{\mathcal U}(C),\operatorname{Pro}_{\mathcal V}(C,\mathcal U)$ & I 9.10.2\\
$\operatorname{Quot}(\mathcal U),Q_X$ & IV 9.8.3 a）\\
$\operatorname{Res}(h)$ & IV 11.1.4\\
$\mathrm{(Topfin)}$ & IV 9.1.12\\
$\operatorname{Top}(X/r)$ & IV 9.8.3 h）\\
$\operatorname{Top}(f)$ & IV 4.9.1, 4.1.1\\
$\operatorname{Top}(X)$ & IV 4.1, 2.1, 4.2\\
$\operatorname{TOP}(f)$ & IV 4.1.3\\
$\operatorname{TOP}(X)$ & IV 2.5, 4.1\\
$u^*$ & I 5.0\\
$u_*,u^*,u_!$ & I 5.1, 5.9；IV 3\\
$u^{-1},u^*,u_*$ & IV 13.2.3\\
$u_s,u^*$ & III 1.1.1\\
$\mathcal U\text{-}\mathrm{Ens}$ & I $0$\\
$\mathcal U,\mathcal V$，宇宙 & I $0$, 11\\
$(\mathcal V,\mathcal U\text{-}\mathrm{Top})$ & IV 3.3\\
$\mathscr{V}\!ois(p)$ & IV 6.8\\
$X_{\complement U}$ & IV 9.3.2\\
$(X_0/U)$ & IV 5.2\\
\end{longtable}
